\begin{abstract}
We evaluate a fixed dictionary-based slang normalization policy on two Indonesian sentiment datasets using separately trained word and character TF-IDF support vector machines and IndoBERT. Shared audited splits, paired uncertainty estimates, direct representation contrasts, and fixed-model input switches distinguish complete pipeline comparisons from inference-only substitution. On SmSA, the normalized-minus-original IndoBERT pipeline effect averaged +0.0202 macro-F1 across three seeds, with a nominal 95\% paired group-bootstrap interval of [0.0037, 0.0381]. Its lower bound did not exceed the investigator-selected +0.01 practical-gain threshold. The other five pipeline intervals and all four direct-contrast intervals included zero. Fixed-model substitution both corrected and introduced errors. These results characterize the tested datasets, models, and dictionary policy rather than establishing a general normalization rule.
\end{abstract}

\begin{IEEEkeywords}
Indonesian sentiment analysis, slang normalization, text preprocessing, IndoBERT, character n-grams
\end{IEEEkeywords}
